% Drafts of the appendix listings that change for sca_v2 (Appendix app:code). Algorithm alg:calib_full does not
% change (calibration is the same; SCAProcessor.calibrate uses the bridge-free local branch), so only its
% neighbour alg:sample_full and the progressive listing alg:prog are given. Each line was checked against
% sca/attn.py (SCAProcessor.sca, global_branch, pool_tokens, coarse_rope, SCAContext.window_now / pi_table_now /
% table_interp / bridge_now, bridge_index, sca_dedup_forward, flex_layout, local_branch_flex, mix_branches,
% PiTable.__call__) and scripts/sample_progressive.py (contexts built with sca/dedup.py context_kwargs).
% New equation labels: eq:sched (window schedule and ratio), eq:bridge (bridge set).

% ---------------------------------------------------------------- replaces the alg:sample_full minipage
\begin{minipage}[t]{0.48\linewidth}
    \hrule height 0.6pt\vspace{2pt}
    \captionof{algorithm}{\textbf{Sampling}, in full}
    \label{alg:sample_full}
    \vspace{-4pt}
\begin{lstlisting}[style=scacode, emph={rope,sdpa,cat,clip,log2,interp,top_k,tile_major,block_mask,flex_attention,inverse}]
# Input: q, k, v before RoPE (T text, then H x W
#   image), block l, noise level sigma
# Output: O for all T + HW queries; sca/attn.py
# Per image: s_H, s_W = H // n0, W // n0; pi_s is
#   the table entry nearest max(s_H, s_W) in log
# Per step, all 57 x 24 heads, ends held:
w = n0                        # Eq. (*@\ref{eq:sched}@*)
for sig_k, w_k in ((.88, 128), (.86, 112),
                   (.84, 96), (.82, 80)):
  if sigma > sig_k: w = w_k; break
if w == n0: pi = interp(sigma, centres, pi_s)
else:                 # ratio r, log-linear
  r = max(H, W) / w; u = clip(log2(r / 2), 0, 1)
  pi = (1-u) * interp(sigma, centres, table[2]) \
     + u * interp(sigma, centres, table[4])
kb = k_img[s_H//2::s_H, s_W//2::s_W]  # (*@$\kappa(b)$@*)
vb = v_img[s_H//2::s_H, s_W//2::s_W]
cp = p // (s_H, s_W)          # (*@$c(p_i)$@*); text at 0
Q = cat(rope(q_txt, 0), rope(q_img, cp))
K = cat(rope(k_txt, 0), rope(kb, b))
O_glob = sdpa(Q, K, cat(v_txt, vb))   # Eq. (*@\ref{eq:glob}@*)
if sigma > sigma_g:                   # 0.89, Eq. (*@\ref{eq:mix}@*)
  return O_glob               # opening phase
# bridge (Eq. (*@\ref{eq:bridge}@*)), per block and step:
d = ((k_img - k_img.mean(0))**2).sum(heads, ch)
B = top_k(d, 16) per cell, 8 x 8 cells  # 1,024
kr = rope(k_img, p)           # true positions
kl = cat(rope(k_txt, 0), kr[B], kr)
vl = cat(v_txt, v_img[B], v_img)  # no dedup
P = T + len(B)                # prefix, seen by all
w = (min(w, H), min(w, W))    # window sides
a = clip(t*(p//t) + t//2 - w//2, 0, (H, W) - w)
keep(i, j) = j < P or a_i <= p_j < a_i + w  # Eq. (*@\ref{eq:loc}@*)
perm = tile_major(H, W)       # 8 x 16 query blocks
mask = block_mask(keep, perm) # 8 x 8 key blocks
O_loc = flex_attention(rope(q_img, p)[perm],
    kl, vl, mask)[inverse(perm)]
O_img = pi[l] * O_loc + (1 - pi[l]) * O_glob[T:]
return cat(O_glob[:T], O_img) # fp32 mix, Eq. (*@\ref{eq:mix}@*)
\end{lstlisting}
\end{minipage}

% ---------------------------------------------------------------- replaces the alg:prog minipage
% Only the use_sca line changes (bridge on; the window schedule is passed but never fires, because every
% stage re-enters below sigma = 0.82).
\vspace{10pt}
\noindent\begin{minipage}{\linewidth}
    \hrule height 0.6pt\vspace{2pt}
    \captionof{algorithm}{\textbf{Progressive \sca{}}, HRDiT's schedule with \sca{} attention (Section~\ref{sec:progressive})}
    \label{alg:prog}
    \vspace{-4pt}
\begin{lstlisting}[style=scacode, emph={shift,linspace,noise,flux,decode,encode,bicubic,sharpen,lowpass,up,momentum,use_sca}]
# Input: prompt, seed; FLUX.1-dev in bf16; SCA tables at scales 2 and 4
# Output: (*@$4096^2$@*) image; scripts/sample_progressive.py
sig = shift(linspace(1, 1/30, 30), mu=1.15)     # one 30-step schedule serves all stages, as in HRDiT
x = noise((*@$1024^2$@*), seed)                         # stage 1: stock dense attention, native RoPE
for i in range(30):                             # Euler; ref keeps each step's x0 estimate
  v = flux(x, sig[i], g=3.5); ref[i] = x - sig[i] * v; x = x + (sig[i+1] - sig[i]) * v
for res, k, g, s_c, a in ((2048, 13, 4.5, 2, 1.0), (4096, 20, 6.0, 4, 0.25)):
  img = sharpen(bicubic(decode(x), res))        # 2 img - blur(img), blur width ~ U(0.1, 2.0)
  x = sig[k] * noise + (1 - sig[k]) * encode(img)   # re-enter at sigma 0.81, then 0.61
  use_sca(table[s_c], n0=64, t=8, sigma_g=0.89, bridge=(8, 16))  # replaces NTK RoPE, attention scale,
                                                # pruning, alignment; w = n0, as sig[k] < 0.82 (Eq. (*@\ref{eq:sched}@*))
  N = 30 - k                                    # 17, then 10 steps; sig[k] < sigma_g: no opening phase
  for n, i in enumerate(range(k, 30)):
    v = flux(x, sig[i], g); x0 = x - sig[i] * v; r, ref[i] = ref[i], x0    # SCA in every block
    x0 = x0 + a * (N - n) / N * (lowpass(up(r)) - lowpass(x0))  # Butterworth, order 4, cut-off 0.2
    u = momentum((x - x0) / sig[i], 0.5 * (N - n) / N)           # HRDiT's velocity momentum
    x = x + (sig[i+1] - sig[i]) * u
return decode(x)
\end{lstlisting}
\end{minipage}
